\section{Análisis de no linealidades}

Las no linealidades de tipo zona muerta y relé son tipos comunes de comportamientos no lineales sin memoria que se encuentran frecuentemente en sistemas físicos. La zona muerta se caracteriza por un rango de valores de entrada dentro del cual la salida del sistema permanece en cero, ignorando la entrada hasta que estas superan ciertos umbrales positivos o negativos. Esto modela fenómenos como la insensibilidad de sensores y el efecto de atascamiento y deslizamiento.

La zona muerta es un fenómeno sin memoria ya que la salida en cualquier instante depende únicamente de la entrada actual, sin considerar la historia de la misma. Esta no linealidad influye profundamente en el desempeño del sistema que debe ser considerado al momento de realizar el modelado, análisis y diseño de sistemas de control.

\subsection{Zona Muerta}


\subsection{Fricción}